\section{Arithmetic obstructions to finite-memory saturation}
\label{sec:r80-arithmetic-saturation}

The convergence in Corollary~\ref{cor:r79-history-limit} need not
terminate at any finite memory.  We give an exact criterion for
termination in a two-model family with rational examples.  Its search
ordering is classical: the index rule and the distinction between
cyclic and acyclic search sequences belong to the search-game
literature \cite{Matula1964,ClarksonLinGlazebrook2023,BuiLidbetterLin2023}.  The issue here is
the implementation of an optimal law by an autonomous stochastic
controller with a counted pure start.  A mixture of finitely many
infinite search sequences need not have a finite positive-memory
implementation.  We prove that distinction, its arithmetic converse,
and a quantitative finite-memory approximation below.

Fix $r,s\in(0,1)$ and write
\[
 q_a=1-r,\qquad q_b=1-s,\qquad
 u=-\log r,\qquad v=-\log s.
\]
There are two scalar models, two commands $a,b$, one failure report,
and compulsory stopping at success.  In model $0$, command $a$
succeeds with probability $q_a$ and $b$ never succeeds.  In model $1$,
command $b$ succeeds with probability $q_b$ and $a$ never succeeds.
The deterministic call costs are arbitrary $c_a,c_b>0$, respectively,
in either model.  Put $\tau_a=c_a/q_a$ and $\tau_b=c_b/q_b$,
the known-model optimal costs.  The model is selected once and remains hidden.
Thus the instrument is feasible, has positive bounded costs, and
has the same observation and initialization conventions as
Section~\ref{sec:r79-quantum-frontier}.  We use $V_M$ and
$V_{\rm hist}$ for its finite-label and unrestricted-history robust
cost values.

\subsection{Search words, the classical exchange rule, and robust duality}

A word $w=(w_0,w_1,\ldots)\in\{a,b\}^{\mathbb N_0}$ specifies what
to do on the hypothetical path on which every call fails.  Let
$i_t,j_t$ be the numbers of $a,b$ in its first $t$ letters.  The
actual costs of that deterministic policy are
\begin{equation}\label{eq:r80-word-costs}
 C_0(w)=\sum_{t\ge0}c_{w_t}r^{i_t},\qquad
 C_1(w)=\sum_{t\ge0}c_{w_t}s^{j_t}.
\end{equation}
These nonnegative sums also define infinite costs without a
convergence convention.  A randomized causal policy induces a
probability law on such words, by sampling its command kernels
along the hypothetical all-failure path.  Conversely every word
law is implemented by its conditional next-letter probabilities.
These probabilities do not require the hidden model: for a fixed
command prefix the physical likelihood of all its failures is
common to every private realization of that prefix and cancels
from the conditional controller law.  Tonelli's theorem therefore
gives the cost vector as the expectation of
\eqref{eq:r80-word-costs} under that word law.

The following search lemma, including interior least-favorable priors,
equalized model costs and a mixture of at most two consistent tie
rules, specializes the classical two-location search theory.  In
particular these conclusions are covered by Proposition 3 and
Theorem 10 of the arXiv version of
Clarkson--Lin--Glazebrook~\cite{ClarksonLinGlazebrook2023}; the
index rule is attributed to Blackwell, as reported by
Matula~\cite{Matula1964}.  We reproduce
a self-contained proof because the finite-memory obstruction
needs its exact statement about all optimal word laws.  These
search and saddle conclusions are not claimed as new results.

\begin{lemma}[Greedy words and an attained interior preparation saddle]
\label{lem:r80-search-saddle}
For a prior $\pi\in(0,1)$, a finite-cost word minimizes
$\pi C_0+(1-\pi)C_1$ if and only if, at every prefix, it chooses
\begin{equation}\label{eq:r80-greedy-rule}
 \begin{cases}
 a,&(\pi/\tau_a)r^{i_t}>((1-\pi)/\tau_b)s^{j_t},\\
 b,&(\pi/\tau_a)r^{i_t}<((1-\pi)/\tau_b)s^{j_t},
 \end{cases}
\end{equation}
with either order permitted at equality.  Every word obeying this
rule has finite cost in both models.  Its asymptotic frequency of
$a$ is
\begin{equation}\label{eq:r80-greedy-frequency}
 \lim_{t\to\infty}\frac{i_t}{t}=\frac{v}{u+v}.
\end{equation}

The robust history value is attained.  There is an interior prior
$\pi_*\in(0,1)$ whose Bayes value is $V_{\rm hist}$.  At that prior,
a robust optimum is a mixture of at most two greedy words: one
places $a$ first at every tied pair, and the other places $b$ first
at every tied pair.  The resulting risk vector is
$(V_{\rm hist},V_{\rm hist})$.
\end{lemma}
\begin{proof}
First consider two consecutive different letters after a prefix
with counts $(i,j)$.  Exchanging $ab$ with $ba$ leaves the counts
after the pair, and hence every later contribution, unchanged.
The two model-cost differences, with $ab$ first in the subtraction,
are exactly
\begin{equation}\label{eq:r80-exchange-vector}
 (-c_bq_a r^i,c_aq_b s^j).
\end{equation}
The Bayes difference is therefore
$c_ac_b(((1-\pi)/\tau_b)s^j-(\pi/\tau_a)r^i)$.  The greedy rule merges the two
strictly decreasing lists
\[
 \frac\pi{\tau_a},\frac\pi{\tau_a}r,
       \frac\pi{\tau_a}r^2,\ldots,
 \qquad
 \frac{1-\pi}{\tau_b},\frac{1-\pi}{\tau_b}s,
       \frac{1-\pi}{\tau_b}s^2,\ldots
\]
in decreasing order.  If a finite-cost competing word disagrees
with a prescribed greedy letter, the desired letter occurs later:
finite cost in both models requires infinitely many occurrences
of both commands.  Move its next occurrence to the desired
position through the intervening copies of the other letter.
Each adjacent exchange is nonincreasing in Bayes cost.  Indeed
the desired list item is at least the first skipped item, and
all subsequent skipped items are smaller.  A strict violation
of \eqref{eq:r80-greedy-rule} gives a strict improvement.

By finitely many such exchanges one can match any prescribed
finite greedy prefix without increasing the competitor's full
cost.  The greedy prefix cost is nonnegative and no greater than
that transformed full cost.  Letting the prefix length increase
proves optimality of every greedy word once its finite cost is
verified.  The strict-improvement observation proves necessity.
An infinite-cost competitor cannot improve a finite Bayes value.

For finiteness and frequency put
\[
 z_0=\log\frac{\pi\tau_b}{(1-\pi)\tau_a},\qquad
 z_t=z_0-ui_t+vj_t.
\]
The rule subtracts $u$ when $z_t>0$, adds $v$ when $z_t<0$, and
allows either operation at zero.  Consequently, for all $t$,
\begin{equation}\label{eq:r80-log-strip}
 \min(z_0,-u)\le z_t\le\max(z_0,v).
\end{equation}
Since $i_t+j_t=t$, the identity
\[
 i_t=\frac{z_0+vt-z_t}{u+v}
\]
proves \eqref{eq:r80-greedy-frequency}.  More precisely, with
$\kappa=\exp(-uv/(u+v))<1$,
\begin{align*}
 r^{i_t}
 &=\exp\!\left(\frac{u(z_t-z_0)}{u+v}\right)\kappa^t,\\
 s^{j_t}
 &=\exp\!\left(\frac{v(z_0-z_t)}{u+v}\right)\kappa^t.
\end{align*}
The strip bounds both factors uniformly, proving convergence of
both cost sums.  These bounds are uniform over all tie choices
at the fixed prior.

We next establish the robust duality without presupposing a
finite optimal word or an infinite-dimensional minimax exchange.
Represent a word law by its cylinder probabilities.  They lie
in $[0,1]$, start with mass one, and satisfy the consistency
equalities and identities for a binary probability tree.  Every
sequence of these laws has a subsequence on which all cylinder
probabilities converge, by successive subsequence selection on
the countable set of cylinders.  The limit again defines a word
law.  Every cost truncated at a fixed number of calls is a
continuous linear function of finitely many cylinders; the full
cost is the supremum of these continuous functions.

Let $\mathcal R\subset\mathbb R^2$ be the upward risk set:
$(x_0,x_1)\in\mathcal R$ when some word law has both costs finite
and no greater than $x_0,x_1$.  It is convex, since word laws can
be mixed.  It is closed: for a convergent sequence of such upper
bounds, take a cylinderwise convergent subsequence of the
associated laws.  Every truncated cost of the limit is bounded
by the limiting coordinate, and monotone convergence gives the
same bound for the full cost.  A fixed full-support command row
makes $\mathcal R$ nonempty.  Thus
$(V_{\rm hist},V_{\rm hist})\in\mathcal R$.

Separate $\mathcal R$ from the disjoint open convex set
$\{(x_0,x_1):x_0<V_{\rm hist},\ x_1<V_{\rm hist}\}$.
A separating linear functional has nonnegative coefficients:
the open set is unbounded in each negative coordinate, while
$\mathcal R$ is upward closed.  Normalize the nonzero
coefficients to $(\pi,1-\pi)$.  The separator and the displayed
boundary point give
\begin{equation}\label{eq:r80-risk-separation}
 \inf_{(x_0,x_1)\in\mathcal R}
          \{\pi x_0+(1-\pi)x_1\}=V_{\rm hist}.
\end{equation}

This prior is interior.  Under model $0$, every proper finite-cost
policy has expected number $1/q_a$ of $a$ calls: summing their
success probabilities gives total success probability one.
Likewise model $1$ requires expected number $1/q_b$ of $b$ calls.
The corresponding cost lower bounds are $\tau_a$ and $\tau_b$.
In fact both inequalities are strict for every law with finite
cost in both models.  Properness in model $1$ requires positive
probability of a $b$ call at some finite commanded prefix.
Every failure in such a prefix has strictly positive probability
in model $0$ as well.  Thus model $0$ incurs a strictly positive
additional expected cost for wrong-command $b$ calls.  The
argument with the models exchanged proves strictness for model
$1$.  Apply this to an attaining robust law: it gives
$V_{\rm hist}>\max(\tau_a,\tau_b)$.

At either endpoint prior the infimum over $\mathcal R$ is the
corresponding $\tau_a$ or $\tau_b$: a full-support memoryless row
can approach the appropriate pure command, with finite cost in
the other model at every stage.  Thus an endpoint cannot satisfy
\eqref{eq:r80-risk-separation}.  Write the resulting interior
prior as $\pi_*$.  We record also a bound used later in the
normalized case $c_a=q_a,c_b=q_b$.  If the first-command
probability of $a$ is $p$, the two costs are at least
$1+q_b(1-p)$ and $1+q_ap$.  Consequently
\begin{equation}\label{eq:r80-bayes-first-call}
 \inf_{\mathcal R}\{\pi x_0+(1-\pi)x_1\}
 \ge 1+\min\{\pi q_b,(1-\pi)q_a\}
 \quad(c_a=q_a,\ c_b=q_b).
\end{equation}

At an interior prior the proved greedy optimum is finite in both
models, so its Bayes value equals the infimum over $\mathcal R$.
An attaining robust law must be supported on Bayes-optimal words:
each word's Bayes cost is at least this value and its expectation
equals the value.  Both of its model costs equal
$V_{\rm hist}$, since both are at most that number and the two
prior weights are positive.

For completeness, tied words reduce to two extreme risk vectors.
The merged numerical lists have a fixed order between unequal
items.  Every tie consists of exactly one item from each list;
its two letters are consecutive.  Changing their order affects
no later count.  Formula~\eqref{eq:r80-exchange-vector} lowers
$C_0$ and raises $C_1$ when the order is $ab$.  At every tie,
$\pi_*r^i/\tau_a=(1-\pi_*)s^j/\tau_b$, so all these cost differences lie on
the same Bayes-level line.  The differences are absolutely
summable, since tied indices are distinct and both geometric
series converge.  Consequently every greedy risk vector lies
on the line segment joining the word with $ab$ at every tie
and the word with $ba$ at every tie.  The same is true of any
randomized greedy law.  The robust risk vector on that segment
is obtained by mixing its two endpoints.  This proves the last
assertion.
\end{proof}

\subsection{The finite-state obstruction}

\begin{lemma}[Deterministic tails of finite stochastic controllers]
\label{lem:r80-finite-word-rigidity}
Suppose an autonomous controller with finitely many labels and
fixed command and update rows, on one perpetually repeated
nonterminal report, emits one prescribed infinite command suffix
with probability one after some fixed finite time.  That suffix is
ultimately periodic.  The conclusion allows arbitrary stochastic
updates and command rows, and does not assume that the retained
label is determined by the visible history.
\end{lemma}
\begin{proof}
Let $S_t$ be the nonempty set of labels of positive probability
after $t$ hypothetical failures.  Once the command suffix is
deterministic, every label in $S_t$ must have the same pure
command row, equal to its prescribed next symbol.  Otherwise a
label of positive probability would emit a different symbol
with positive probability.  Given that symbol, the next support
is the union of the supports of the corresponding update rows
from labels in $S_t$.  Thus $S_{t+1}$, and its next prescribed
symbol, are determined by $S_t$.  There are only finitely many
nonempty subsets of the finite label set.  Once one support
recurs after the suffix begins, its entire future support and
symbol sequence recur.  This proves ultimate periodicity.
The argument uses supports, not the more restrictive assumption
of deterministic internal transitions.
\end{proof}

\begin{theorem}[Arithmetic dichotomy and exponentially accurate finite memory]
\label{thm:r80-arithmetic-saturation}
For the two-model family above,
\[
 V_1=\tau_a+\tau_b,
 \qquad
 \bigl(\exists M<\infty:\ V_M=V_{\rm hist}\bigr)
       \quad\Longleftrightarrow\quad
                  \frac{\log r}{\log s}\in\mathbb Q.
\]
If $\log r/\log s$ is irrational, then every finite budget has
an attained strict gap
\begin{equation}\label{eq:r80-strict-all-memory}
 V_M>V_{\rm hist},\qquad M=1,2,\ldots,
 \qquad V_M\downarrow V_{\rm hist}.
\end{equation}
This statement ranges over all randomized autonomous controllers
with a prescribed pure initial label, not only deterministic
cyclic command schedules.

For the normalized costs $c_a=q_a,c_b=q_b$, there is an explicit
upper estimate for both arithmetic cases.  Put
\begin{gather*}
 m=\min(q_a,q_b),\qquad \delta=\frac{m^2}{16},\qquad
 L=\log\frac{1-\delta}{\delta},\\
 A=\exp\!\left(\frac{uv+\max(u,v)L}{u+v}\right),
 \qquad \kappa=\exp\!\left(-\frac{uv}{u+v}\right).
\end{gather*}
Then for every $M\ge2$,
\begin{equation}\label{eq:r80-memory-exponential}
 0\le V_M-V_{\rm hist}
                 \le 2A\kappa^{\lfloor M/2\rfloor}.
\end{equation}
The upper rate is not claimed to be a matching lower rate or an
exact formula for intermediate memory values.  In the rational
case a finite implementation is constructed below, without a
claim that its label count is minimal.
\end{theorem}
\begin{proof}
For a one-label row with command-$a$ probability $p\in(0,1)$ put
$x=pq_a$ and $y=(1-p)q_b$.  Its mean cost per call is
$\tau_a x+\tau_b y$, and its model success probabilities are
$x$ and $y$.  The costs are
$\tau_a+\tau_b y/x$ and $\tau_b+\tau_a x/y$.
If $x\ge y$ the second is at least $\tau_a+\tau_b$; if $x\le y$
the first is.  Equality holds at $x=y$, or
$p=q_b/(q_a+q_b)$.  Endpoint rows are improper in one model.
This proves $V_1=\tau_a+\tau_b$.  In particular the normalized
case has a common one-label fallback of robust cost two.

Take the interior prior $\pi_*$ from
Lemma~\ref{lem:r80-search-saddle}.  If $u/v$ is irrational, there
is at most one tied pair: two distinct equalities
$\pi_*r^i/\tau_a=(1-\pi_*)s^j/\tau_b$ would imply
$u(i-i')=v(j-j')$, contrary to irrationality.  Thus there is one
greedy word, or two which differ only in the order of a single
adjacent pair and have an identical suffix.  Every such suffix
has the irrational frequency $v/(u+v)$ by
\eqref{eq:r80-greedy-frequency}.  It is therefore not ultimately
periodic.

Suppose some finite-label controller attained the robust history
value.  Its planned-word law would be Bayes-optimal at $\pi_*$,
and hence supported on those one or two words.  After their
possible single tied pair its command suffix would be prescribed
with probability one.  Lemma~\ref{lem:r80-finite-word-rigidity}
would make that suffix ultimately periodic, a contradiction.
All failure probabilities in this family are strictly positive,
so the hypothetical-failure support argument introduces no
physical zero-likelihood history.  At each fixed $M$ the optimum
is attained by Theorem~\ref{thm:r79-quantum-frontier}; therefore
the absence of an optimal controller at that $M$ gives the strict
inequality in \eqref{eq:r80-strict-all-memory}, rather than only
nonattainment.  Its decreasing convergence follows either from
Corollary~\ref{cor:r79-history-limit} or from the explicit
construction below.

Suppose instead that $u/v$ is rational.  Write $u=h\ell$,
$v=k\ell$ with positive integers $h,k$ and $\ell>0$.
For either fixed tie rule, the logarithmic state $z_t$ belongs
to the coset $z_0+\ell\mathbb Z$ and to the bounded strip
\eqref{eq:r80-log-strip}.  This intersection is finite.  The
fixed tie rule makes its next state a deterministic function
of its present state, so each of the two extreme greedy words
is ultimately periodic.  An ultimately periodic word has a
finite autonomous realization by one label for each position
in its finite prefix and cycle.

The two-word robust mixture in
Lemma~\ref{lem:r80-search-saddle} also has a finite realization
with a counted pure start.  Use disjoint finite realizations of
the two tails and one initial label.  If their mixing weights
are $\lambda_j$ and their first symbols are $a_j$, the initial
command row assigns probability
$p_a=\sum_{j:a_j=a}\lambda_j$ to symbol $a$.  On its failure,
select tail $j$ with probability $\lambda_j/p_a$ among those
with $a_j=a$; an empty class is never selected.  These are
allowed command-dependent stochastic updates.  Conditional on
the sampled command, its failure probability is independent of
which such tail will be selected, so this constructs exactly
the desired mixture in both models.  Its branch choice is
stored in counted disjoint tail labels.  No freely randomized
initial label or persistent external coin is used.  The finite
implementation proves the rational direction.

For the quantitative approximation, now assume
$c_a=q_a,c_b=q_b$, so $\tau_a=\tau_b=1$.  The constants
will depend only on $r,s$.  A memoryless row with
$p/(1-p)=(q_b/q_a)\sqrt{\pi/(1-\pi)}$ has Bayes cost
\[
 1+2\sqrt{\pi(1-\pi)}.
\]
Together with \eqref{eq:r80-bayes-first-call} at $\pi=1/2$ this
gives, at every least-favorable interior prior,
\[
 1+\frac m2\le V_{\rm hist}
        \le1+2\sqrt{\pi_*(1-\pi_*)}.
\]
Consequently $\pi_*(1-\pi_*)\ge m^2/16$ and
$\delta\le\pi_*\le1-\delta$, so $|z_0|\le L$.
The strip and the exact survival expressions in the proof of
Lemma~\ref{lem:r80-search-saddle} now give, simultaneously for
both extreme words and both models,
\begin{equation}\label{eq:r80-greedy-tail}
 \Pr(\text{first }T\text{ calls all fail})
                       \le A\kappa^T\qquad(T\ge0).
\end{equation}

Follow the robust two-word mixture for at most $T\ge1$ calls;
after failure of its $T$th call, switch to the one-label fallback
of cost two.  The truncated prefix costs at most
$V_{\rm hist}$ in either model, and the added expected cost is
at most $2A\kappa^T$ by \eqref{eq:r80-greedy-tail}.  This is an
actual counted controller with at most $2T$ labels: one pure
initial label, at most $2(T-1)$ disjoint labels for the remaining
prefix positions of the two words, and one fallback label.
The same command-conditioned initial update used above realizes
the mixture.  Depth and branch identity are stored in these
labels.  Thus $V_{2T}\le V_{\rm hist}+2A\kappa^T$.
Taking $T=\lfloor M/2\rfloor$ and allowing unused labels proves
\eqref{eq:r80-memory-exponential}.
\end{proof}

\begin{corollary}[An effective saturation test for rational primitives]
\label{cor:r80-effective-saturation}
Suppose $r,s\in\mathbb Q\cap(0,1)$.  Factor their numerators and
denominators into primes, and let $\nu(r),\nu(s)$ be their signed
prime-exponent vectors on the union of those primes.  The history
optimum has a finite-label realization if and only if these vectors
are positive rational multiples of one another.  This is a finite
exact test on the primitive data; it does not search unboundedly
over label budgets.

In the affirmative case one can compute a rational $t\in(0,1)$
and coprime positive integers $h,k$ with $r=t^h$, $s=t^k$.
For the normalized costs $c_a=q_a,c_b=q_b$, keep $\delta$ from Theorem~\ref{thm:r80-arithmetic-saturation},
and choose any integer $J\ge0$ with
\[
 t^{-J}\ge\frac{1-\delta}{\delta}.
\]
Then the explicitly computable label bound
\begin{equation}\label{eq:r80-rational-label-cap}
 M_0=1+2\bigl(\max(J,h)+\max(J,k)+1\bigr)
\end{equation}
satisfies $V_{M_0}=V_{\rm hist}$.  The history value and an
attaining finite controller are therefore algebraically computable
by the fixed-budget procedure of
Theorem~\ref{thm:r79-quantum-frontier}.  Neither $M_0$ nor the
factorization procedure is asserted to be optimal in its resource
use.
\end{corollary}
\begin{proof}
A rational logarithmic ratio is equivalent to $r^k=s^h$ for some
positive integers $h,k$, and unique factorization makes this
equivalent to $k\nu(r)=h\nu(s)$.  The latter is tested by finitely
many integer cross-products.  Reducing $h,k$ to coprime integers
shows that every coordinate of $\nu(r)$ is divisible by $h$ and
every coordinate of $\nu(s)$ by $k$.  Their common quotient
vector defines the rational number $t$ with $r=t^h$, $s=t^k$.
It lies in $(0,1)$ because $r$ does.  Repeated rational multiplication
finds an integer $J$ satisfying the displayed inequality.

Put $\ell=-\log t$.  The least-favorable logarithmic prior obeys
$|z_0|\le L\le J\ell$.  The entire orbit of either fixed tie rule
lies in its coset $z_0+\ell\mathbb Z$ and in
\[
 [-\max(J,h)\ell,\ \max(J,k)\ell].
\]
There are at most $N=\max(J,h)+\max(J,k)+1$ points of that
coset in this interval.  The deterministic orbit, including its
transient and cycle, therefore uses at most $N$ states.  Each
extreme word and each of its tails has a realization with at
most $N$ labels.  The counted initial-mixture construction in
the preceding proof uses at most $1+2N=M_0$ labels and attains
the robust history value.  The fixed-$M_0$ algebraic synthesis
uses only the rational instrument and costs; it does not require
an exact evaluation of the logarithmic prior.  This proves the
effective conclusion.
\end{proof}

\begin{example}[A rational primitive with no finite saturation]
\label{ex:r80-rational-nonsaturation}
Take $r=1/2$ and $s=1/3$.  The scalar nonterminal multipliers are
$(1/2,1)$ in model $0$ and $(1,1/3)$ in model $1$, and the
command costs are $1/2,2/3$.  All instrument and cost data are
rational.  The ratio $\log2/\log3$ is irrational: a rational
ratio would give $2^n=3^m$ for positive integers $m,n$.
Therefore $V_1=2$ and $V_M>V_{\rm hist}$ at every finite memory,
although \eqref{eq:r80-memory-exponential} supplies exponentially
accurate finite controllers.  Each fixed-$M$ value and an
attaining controller are algebraically computable by
Theorem~\ref{thm:r79-quantum-frontier}.  No transcendence assertion
about $V_{\rm hist}$ is needed or made.
\end{example}

This is an exact saturation classification for the stated
search-type family, not a classification of arbitrary quantum
instruments or general partially observed control.  The exchange
ordering is inherited from classical search theory.  The
finite-support rigidity argument is what excludes every
stochastic finite-label implementation of its irrational optimal
law.  The positive approximation bound and the negative exact
realization statement address different resource questions;
neither replaces the original squared-sign risk frontier or the
independent geometric and Theta construction problems.
